% ============================================================
%  About This Work: SparseRenderFormer
%  RenderFormer Essentials, Main Contribution, Supporting Contribution
%  Internal document. Compile: tectonic about_sparserenderformer.tex
% ============================================================

\documentclass[11pt,letterpaper]{article}

\usepackage[T1]{fontenc}
\usepackage[utf8]{inputenc}
\usepackage{geometry}
\geometry{top=1in, bottom=1in, left=1.25in, right=1.25in}

\usepackage{amsmath, amssymb, amsthm}
\usepackage{booktabs}
\usepackage{tabularx}
\usepackage{enumitem}
\usepackage{xcolor}
\usepackage{hyperref}
\usepackage{microtype}
\emergencystretch=2em
\usepackage{parskip}
\usepackage{titlesec}
\usepackage{mdframed}

\hypersetup{colorlinks=true, linkcolor=blue!60!black,
            citecolor=green!50!black, urlcolor=blue!60!black}

\titleformat{\section}{\large\bfseries}{\thesection.}{0.5em}{}
\titleformat{\subsection}{\normalsize\bfseries}{\thesubsection}{0.5em}{}

\newcommand{\ours}{SparseRenderFormer}
\newcommand{\RF}{RenderFormer}
\newcommand{\step}[1]{\textbf{Step #1.}}
\newcommand{\code}[1]{\texttt{#1}}

\newmdenv[linecolor=blue!40, backgroundcolor=blue!5,
          innerleftmargin=8pt, innerrightmargin=8pt,
          innertopmargin=6pt, innerbottommargin=6pt]{claimbox}

\newmdenv[linecolor=orange!60, backgroundcolor=orange!5,
          innerleftmargin=8pt, innerrightmargin=8pt,
          innertopmargin=6pt, innerbottommargin=6pt]{actionbox}

% ============================================================
\begin{document}

\begin{center}
  {\LARGE\bfseries About This Work: \ours}\\[6pt]
  {\large RenderFormer Essentials $\cdot$ Main Contribution $\cdot$ Supporting Contribution}\\[4pt]
  {\normalsize Victor Miene \quad|\quad Carnegie Mellon University}\\[2pt]
  {\small\color{gray}Internal document. Not for submission. \today}
\end{center}

\hrule
\vspace{10pt}

\tableofcontents
\newpage

% ============================================================
\section*{The Problem in Plain Terms}

RenderFormer works by having every triangle in a 3D scene compare itself
to every other triangle to compute how light bounces between them.
That comparison grows quadratically: double the number of triangles and
the work quadruples.
At 4{,}096 triangles this already pushes the limit of a high-end GPU.
Real 3D assets (rooms, characters, product models) routinely
have 8{,}000 to 32{,}000 triangles or more.
RenderFormer runs out of memory before it can render them.

\ours fixes this by making the comparisons selective.
Instead of every triangle looking at every other triangle, each triangle
looks at three targeted groups: a compressed summary of the whole scene,
a small set of surfaces that are physically likely to matter for its
lighting, and its immediate spatial neighbours.
The total number of comparisons drops from quadratic to roughly linear,
which allows the method to scale to scenes that were previously too large
to render with this approach.

This change affects both inference and fine-tuning, but for the same
underlying reason and with different intent.
At inference, a 32{,}768-triangle scene simply does not fit in GPU memory
under dense attention --- the method cannot run at all.
GP-Sparse makes the forward pass fit, so rendering large scenes becomes
possible.
During fine-tuning the same forward pass is required, so GP-Sparse also
makes training on large scenes tractable.
The intent, however, is inference: cheaper fine-tuning is a side effect,
not a goal.
This distinguishes \ours from NSA, which explicitly targets both training
speed and inference speed as separate objectives across different phases
of an autoregressive model.

% ============================================================
\section{RenderFormer Essentials}

\begin{mdframed}[linecolor=black!30, backgroundcolor=gray!5]
\textbf{What this section covers.}
The components of \RF that \ours directly modifies or depends on.
Read this before the additions below.
\end{mdframed}

\subsection*{Must-Know Bullet Points}

\begin{itemize}

  \item \textbf{Triangle token (13 dimensions).}
    Each triangle is represented as a 13-dimensional feature vector:
    surface normal (3 dims), diffuse albedo (3), specular albedo (3),
    roughness (1), and emission (3).
    Geometric position is not concatenated; it is encoded via 3D~RoPE
    (see next bullet).
    The normal slice is \code{[:, 0:3]} in the token tensor.

  \item \textbf{3D RoPE (not standard positional encoding).}
    Standard positional encoding adds a position vector directly to the
    token embedding before attention.
    Position therefore becomes part of the token content, and the model
    must learn to disentangle geometry from appearance.

    RoPE (Rotary Position Encoding, Su et al.\ RoFormer) works differently.
    It does not touch the token at all.
    Instead, it rotates the query vector $\mathbf{q}_i$ and the key vector
    $\mathbf{k}_j$ in the attention computation using their respective
    positions, so that the dot product $\mathbf{q}_i \cdot \mathbf{k}_j$
    automatically encodes the \emph{relative} position between tokens $i$
    and $j$.
    No position information is stored in the token; it enters only at the
    moment two tokens are compared.

    \RF extends this to three spatial dimensions.
    Each triangle has three vertex positions
    $\mathbf{v}^{(1)}, \mathbf{v}^{(2)}, \mathbf{v}^{(3)} \in \mathbb{R}^3$.
    These are used to rotate $\mathbf{q}$ and $\mathbf{k}$ so that the
    attention logit between triangles $i$ and $j$ reflects their relative
    3D geometry.
    The result is translation invariant by construction.
    There are no learned positional embeddings.

    In GP-Sparse, 3D RoPE is applied per branch:
    the coarse branch uses the centroid of each BVH block as the position
    for the compressed block token; the selected and local branches use
    each triangle's own vertex positions.
    This means geometric relationships are preserved at block resolution
    in the coarse branch and at triangle resolution in the other two.

  \item \textbf{View-independent stage (the bottleneck we fix).}
    Twelve standard self-attention layers propagate triangle-to-triangle
    light transport.
    Complexity is $O(N^2)$ in time and memory.
    At $N = 4{,}096$ this fits on one A100-80GB; at $N = 32{,}768$ it
    requires $64\times$ more memory and is infeasible.
    \ours replaces this stage with three-branch sparse attention.

  \item \textbf{View-dependent stage.}
    Six cross-attention layers map ray-bundle tokens to triangle features
    to produce HDR pixel outputs.
    Complexity is $O(R \cdot N)$ for $R$ ray bundles.
    \ours reduces this to $O(R \cdot k_r)$ via frustum culling.

  \item \textbf{Pixel decoder.}
    A Swin-Large backbone converts ray-bundle features to HDR pixel
    patches. \ours leaves this untouched.

  \item \textbf{Training.}
    205M parameters trained on 16M HDR Objaverse renders for 8 days on
    8$\times$A100-80GB.
    \ours fine-tunes from this checkpoint rather than retraining from
    scratch; only the attention projection matrices in the view-independent
    stage are re-initialised.

    \textbf{Why fine-tuning is sufficient, and why we do not need heavy GPU.}
    The 8-day run taught \RF everything from scratch: how materials look,
    how light bounces, what geometry means.
    That knowledge is distributed across the token embeddings, the
    view-dependent cross-attention layers, and the Swin-L pixel decoder.
    None of those weights change.
    The only weights that change are the attention projections in the
    view-independent stage, because those are the ones whose access
    pattern changes from dense to sparse.
    The model does not need to relearn physics or appearance; it only
    needs to adapt to attending selectively rather than exhaustively.

    The architectural change also makes large-scale fine-tuning tractable
    in a second way.
    Dense attention at $N = 32{,}768$ requires $64\times$ more memory than
    at $N = 4{,}096$; a forward pass at that scale is infeasible with the
    original \RF architecture.
    GP-Sparse reduces the view-independent stage from $O(N^2)$ to
    effectively $O(N)$, so scenes that were previously unrunnable now fit
    in GPU memory.
    The architectural change is therefore not a training cost but what
    makes fine-tuning at target scale possible at all.

    \textbf{The primary benefit is at inference, not training.}
    The goal of GP-Sparse is to render scenes with far more triangles
    than \RF can handle at all.
    Training is affected only because GP-Sparse makes a forward pass at
    large $N$ feasible; without it, fine-tuning on large scenes would
    be impossible.
    But cheaper or faster training is not the objective.
    The objective is inference on scenes that were previously unrunnable.

  \item \textbf{What we change: GP-Sparse (Geometry-Physics Sparse Attention).}
    The view-independent self-attention is replaced with GP-Sparse,
    a three-branch module whose block structure follows the scene BVH.
    GP-Sparse is the main contribution of this paper.
    The two physics-grounded additions described in Sections~\ref{sec:main}
    and~\ref{sec:supporting} are specific modifications within GP-Sparse.
    \begin{itemize}
      \item \textbf{Coarse branch.} Each triangle attends to area-weighted
        compressed summaries of \emph{parent-level} BVH blocks (one level
        above the leaf nodes). Using the parent level produces tokens that
        each cover at least $2\times$ more triangles than a leaf, giving
        the branch a genuinely coarse receptive field. Captures low-frequency
        radiance such as ambient light and large colour transfers.
      \item \textbf{Selected branch.} Each triangle selects the top-$k$
        \emph{leaf-level} BVH blocks using a lightweight block scorer:
        \[
          I_{i,\beta} = \sum_{j=1}^{H^I} w^I_{i,j}
                        \cdot \mathrm{ReLU}\!\left(q^I_{i,j}
                        \cdot \bar{k}^I_\beta\right),
        \]
        where $q^I_{i,j}$ derives from triangle $i$'s token,
        $\bar{k}^I_\beta$ derives from the area-weighted pooled token of
        block $\beta$, and $H^I$ is the number of scorer heads (e.g.\
        $H^I = 4$). ReLU activation suppresses negative scores without
        normalisation, which is hardware-efficient and sparse by
        construction. In Stage~1 the scorer is calibrated via KL warm-up;
        in Stage~2 the top-$k$ selection uses a straight-through estimator
        (STE) so rendering-loss gradients reach the scorer and main model
        jointly (see the training procedure below). Captures direct
        illumination, one-bounce specular, and shadow casters.
      \item \textbf{Local branch.} Each triangle attends to its 64
        spatially nearest neighbours in 3D, with back-facing neighbours
        suppressed. Captures contact shadows and near-field colour bleeding.
    \end{itemize}
    The outputs of all three branches are combined by per-head learned
    gates $\lambda_1, \lambda_2, \lambda_3$.
    Both contributions introduce zero new parameters.

\end{itemize}

% ------------------------------------------------------------
\subsection*{Input Encoding: How a Scene Becomes Tokens}

\textbf{Step 1: Each triangle becomes a 13-dimensional feature vector.}
Surface attributes are packed in a fixed order; geometry is handled
separately (see Step~2).

\begin{center}\small
\begin{tabular}{@{}llc@{}}
\toprule
Slice & Attribute & Dims \\
\midrule
\code{[:, 0:3]}   & Surface normal    & 3 \\
\code{[:, 3:6]}   & Diffuse albedo    & 3 \\
\code{[:, 6:9]}   & Specular albedo   & 3 \\
\code{[:, 9]}     & Roughness         & 1 \\
\code{[:, 10:13]} & Emission          & 3 \\
\bottomrule
\end{tabular}
\end{center}

\textbf{Step 2: Geometric position enters via 3D RoPE, not the token.}
Standard positional encoding adds a position vector to each token before
attention, forcing the model to disentangle geometry from appearance.
3D RoPE never touches the token.
Instead, it rotates the query $\mathbf{q}_i$ and key $\mathbf{k}_j$ at
attention time using each triangle's three vertex positions, so the dot
product $\mathbf{q}_i \cdot \mathbf{k}_j$ encodes the relative 3D
geometry between triangles $i$ and $j$ automatically.
The result is translation-invariant by construction and requires no
learned positional embeddings.

In GP-Sparse, 3D RoPE is applied per branch:
the coarse branch uses the centroid of each BVH block as the position for
the compressed block token (block-resolution geometry); the selected and
local branches use each triangle's own vertex positions (triangle-resolution
geometry).

\textbf{Step 3: Tokens flow through the two-stage pipeline.}

\begin{center}\small
\begin{tabular}{@{}cl@{}}
\toprule
Stage & Operation \\
\midrule
Input & Triangle tokens $(N \times 13)$ \\
View-independent & 12$\times$ self-attention layers (3D RoPE on $\mathbf{q}$, $\mathbf{k}$) \\
                 & \textit{Replaced by GP-Sparse in \ours} \\
View-dependent   & 6$\times$ cross-attention: ray-bundle queries attend to triangle keys/values \\
Pixel decoder    & Swin-Large $\to$ HDR pixel patches \\
Output           & HDR image \\
\bottomrule
\end{tabular}
\end{center}

The 13-dimensional token format is preserved throughout the view-independent
stage; nothing is compressed or re-encoded between layers.
After the 12 self-attention layers, each token carries the same dimensions
but is now contextualised: it has gathered information from other triangles
via attention.
The view-dependent stage then changes modality (mesh tokens to pixels) via
cross-attention, not via a bottleneck.

\subsection*{Block Scorer and Training Procedure}

GP-Sparse uses a two-stage continued training procedure, adapted from DSA
(DeepSeek-AI, 2025), to initialise the block scorer before sparse routing
is imposed on a dense RenderFormer checkpoint.

\textbf{Stage 1: Dense Warm-up.}
All GP-Sparse parameters are frozen except the block scorer.
Full dense attention runs normally.
The scorer is trained with a KL divergence loss against the dense attention
distribution:
\begin{equation}
  \mathcal{L}^I = \sum_i \mathbb{D}_{\mathrm{KL}}\!\left(
    p_{i,:} \;\Big\|\; \mathrm{Softmax}(I_{i,:}) \right),
  \label{eq:scorer_warm_up}
\end{equation}
where $p_{i,:}$ is the L1-normalised sum of the dense attention weights for
triangle $i$ across all heads.
This stage runs for a small number of steps and teaches the scorer to predict
which BVH blocks the dense model finds important before any sparse routing is
imposed.
Without this warm-up, sparse routing starts from a random initialisation and
the indirect gradient from the rendering loss is too weak to align the scorer
reliably.

\textbf{Stage 2: Sparse Fine-Tuning.}
Sparse block selection and the normal-coherent mask are activated; all
parameters are unfrozen.
Both discrete operations use straight-through estimators (STE): the
forward pass uses the hard top-$k$ selection and the hard normal mask;
the backward pass passes gradients through the scorer logits and the
sigmoid soft-mask respectively.
This allows the rendering loss to flow end-to-end into both the scorer
and the triangle normal representations.
The auxiliary $\mathcal{L}^I$ continues as a calibration signal, weighted
by $\alpha = 0.01$:
$\mathcal{L} = \mathcal{L}^{\text{render}} + \alpha\,\mathcal{L}^I$.

\subsection*{Where GP-Sparse Comes From}

GP-Sparse is adapted from \textbf{NSA (Native Sparse Attention)},
introduced by Yuan et al.\ (arXiv:2502.11089) for language modelling.
NSA decomposes attention into three hardware-aligned branches:
compressed token blocks, dynamically selected fine blocks, and a local
sliding window.
It achieves near-lossless perplexity at $O(N)$ cost on sequences of up
to 64k tokens.

GP-Sparse takes the three-branch structure of NSA and replaces every
language-specific component with a rendering-equation-grounded equivalent:

\begin{center}\small
\begin{tabular}{@{}ll@{}}
\toprule
NSA & GP-Sparse \\
\midrule
Sequential block boundaries     & BVH leaf node boundaries \\
Uniform pooling                 & Area-weighted pooling (rendering integral) \\
1D causal RoPE                  & 3D RoPE on triangle vertex positions \\
Local sliding window            & $k$-NN in 3D space \\
No neighbourhood mask           & Normal-coherent mask (hemisphere constraint) \\
\midrule
Single block depth for all branches
  & Parent level (coarse branch) $+$ leaf level (selected branch) \\
Compressed attn as routing signal
  & Dedicated block scorer, KL warm-up, STE end-to-end gradient \\
Overlapping stride ($d < l$)
  & Boundary fringe ($k_{\text{fringe}}$ triangles from neighbouring leaves) \\
\bottomrule
\end{tabular}
\end{center}

NSA and all prior sparse attention methods operate on one-dimensional
token sequences with no spatial geometry.
Their neighbourhood masks carry no physical meaning.
GP-Sparse is the first method to apply a geometric visibility primitive
(back-face culling) as a hard attention mask, which is precisely what
Addition~1 (Section~\ref{sec:main}) contributes.
Addition~2 (Section~\ref{sec:supporting}) grounds the coarse branch
pooling step in the rendering integral, which NSA does not do.

% ============================================================
\section*{Things to Note When Reading the Literature}

\begin{mdframed}[linecolor=black!30, backgroundcolor=gray!5]
Reading comprehension notes for the key papers in this area.
What to look for and how each paper relates to this work.
\end{mdframed}

\paragraph{NSA (Yuan et al., arXiv:2502.11089).}
This is the direct structural source of GP-Sparse.
When reading it, note that the three-branch split (compressed blocks,
selected blocks, local window) is the core idea, and that every branch
is designed around one-dimensional token sequences with fixed sequential
block boundaries.
There is no notion of geometry, physical plausibility, or neighbourhood
semantics.
The local window is simply a fixed sliding window over token indices.
Understanding this gap is essential: GP-Sparse replaces every
sequence-specific component with a geometry- or physics-grounded
equivalent.

\paragraph{RenderFormer (Zeng et al., SIGGRAPH 2025).}
This is the base model.
When reading it, pay close attention to: (1) the triangle token format
and what each dimension encodes; (2) how 3D RoPE encodes position in
the attention logits rather than the token; (3) the separation between
the view-independent and view-dependent stages and why that separation
matters for reuse across viewpoints; (4) the training data pipeline
from Objaverse.
Everything in GP-Sparse either modifies or depends on one of these four
components.

\paragraph{The rendering equation (Kajiya, 1986).}
This paper defines the physical law that both additions are grounded in.
Note the hemisphere integral: radiance arriving at a surface point
integrates only over directions in the upper hemisphere
($\mathbf{n} \cdot \omega > 0$).
This is the physical justification for the normal-coherent mask
($\tau = 0.5$).
Also note that the integral weights a surface patch by its projected
area ($\cos\theta \, dA$), which is the justification for area-weighted
pooling in the coarse branch.

\paragraph{FlashAttention (Dao et al., 2022, 2023).}
When reading this, note that FlashAttention reduces memory \emph{I/O}
by tiling the softmax computation in SRAM.
It does not reduce the number of attention pairs computed; complexity
remains $O(N^2)$.
GP-Sparse and FlashAttention address different problems: GP-Sparse
reduces which pairs are computed; FlashAttention makes computing a
given set of pairs more memory-efficient.
The two are complementary.

\paragraph{BigBird, Longformer, Reformer.}
These are earlier sparse attention methods.
When reading them, note that sparsity is defined by sequence position
(local windows, global tokens, random hashing).
None use geometry to define neighbourhoods and none operate on 3D
structure.
They are worth knowing for context but NSA is the relevant comparison
for GP-Sparse because NSA is also hardware-aligned and multi-branch.

\paragraph{NeRF, 3DGS, DiffusionRenderer.}
When reading these, note what they do \emph{not} do: none model
explicit inter-surface light transport.
NeRF and 3DGS encode appearance implicitly and bake in the lighting
at training time; changing a light source requires retraining or
re-optimising.
DiffusionRenderer uses G-buffers from a rasteriser and cannot operate
on arbitrary mesh geometry directly.
\RF is different because it propagates light between explicit triangle
tokens, enabling relighting without per-scene baking.
Keep this distinction clear when reading any neural rendering paper.

\paragraph{EquiformerV2 and SE(3)-equivariant networks.}
When reading these, note what equivariance gives you: a model that
produces geometrically consistent outputs regardless of how the scene
is rotated in 3D.
\RF and \ours do not have this property.
3D~RoPE gives translation invariance but not rotation invariance;
\RF compensates with rotation augmentation during training.
This is a known limitation of the current design, not an oversight.

\paragraph{BVH literature (surface area heuristic, SAH).}
The coarse branch and selected branch both rely on BVH structure.
When reading BVH papers, note how SAH splitting works: it minimises
expected ray-intersection cost by splitting where the product of
surface area and primitive count is minimised.
A side effect is that SAH tends to produce leaves of similar-sized
triangles, which means area weights within a leaf may be nearly
uniform on well-processed meshes.
This is one of the listed disadvantages of Addition~2.

% ============================================================
\section*{NSA: Foundations This Work Builds On}

GP-Sparse is structurally derived from NSA.
Understanding what NSA is built on helps clarify what GP-Sparse inherits
and what it replaces.

\subsection*{1. Architectural Foundations}

NSA is built on the standard Transformer and its vanilla attention
mechanism.
Three modern architectural choices shape the design:

\begin{itemize}[nosep]
  \item \textbf{Grouped-Query Attention (GQA) and Multi-Query Attention
    (MQA).} NSA is designed to be compatible with these architectures,
    which share KV-caches across multiple query heads to reduce memory
    bandwidth.
    GP-Sparse does not use GQA (each triangle token has its own head),
    so this compatibility constraint does not carry over.
  \item \textbf{Mixture-of-Experts (MoE).} NSA experiments use a backbone
    combining GQA with DeepSeekMoE, totalling 27B parameters with only
    3B active per token.
    GP-Sparse uses the standard \RF attention projections without any
    expert routing.
  \item \textbf{YaRN.} NSA uses YaRN to extend context from 8k to 32k
    tokens during fine-tuning.
    GP-Sparse handles long context differently: the BVH partitions the
    sequence geometrically so attention cost grows linearly with $N$
    from the start, without any context-extension fine-tuning.
\end{itemize}

\subsection*{2. Hardware and Implementation Foundations}

The ``native'' and ``hardware-aligned'' aspects of NSA rest on two
implementation decisions:

\begin{itemize}[nosep]
  \item \textbf{FlashAttention principles.} NSA adopts blockwise memory
    access and contiguous computation, maximising Tensor Core utilisation.
    GP-Sparse inherits this: BVH leaf blocks map directly to contiguous
    memory, making blockwise kernels applicable to the selected branch.
    The local branch (kNN) is less contiguous and is the main
    implementation risk.
  \item \textbf{Triton kernels.} NSA implements its sparse attention
    in Triton for tiled computation.
    GP-Sparse can port the NSA Triton kernel (available in the
    DeepSeek-V3 repo) with minimal changes; the block boundaries change
    from sequential strides to BVH leaf offsets.
\end{itemize}

\subsection*{3. Conceptual and Empirical Inspirations}

Three observations drove NSA's design choices:

\begin{itemize}[nosep]
  \item \textbf{Inherent sparsity.} Softmax attention is naturally sparse:
    only a small subset of query-key pairs contributes significantly to
    the output.
    GP-Sparse exploits the same property but uses geometry (BVH
    containment, surface normal alignment) to \emph{predict} which pairs
    matter, rather than learning it purely from data.
  \item \textbf{Blockwise clustering pattern.} Attention maps from a
    pretrained full-attention model show that nearby keys share similar
    importance scores.
    This directly justifies blockwise selection over token-granular
    selection: contiguous blocks are hardware-efficient and recover the
    same important regions.
    GP-Sparse grounds this observation in physics: triangles within the
    same BVH leaf share a local light-transport context, which is a
    geometry-based explanation for why block selection works in the
    rendering domain.
  \item \textbf{Critique of post-hoc methods.} NSA identifies that H2O,
    MInference, Quest, and ClusterKV fail because they apply sparsity
    only at inference time or use non-differentiable operations that
    prevent the model from learning optimal sparse patterns.
    GP-Sparse addresses both issues: the two-stage training procedure
    (Section~\ref{sec:main}) trains with sparsity from the start, and the
    straight-through estimators on the top-$k$ selection and normal mask
    maintain full differentiability throughout fine-tuning.
\end{itemize}

In summary, NSA evolves the Transformer by integrating algorithmic
sparsity with hardware-centric implementation for long-context
training and inference.
GP-Sparse inherits the three-branch structure and blockwise kernel design,
and replaces every sequence-specific component with a rendering-equation
equivalent.

\subsection*{4. Optimization Goals: NSA vs.\ \ours}

NSA targets two distinct bottlenecks depending on phase, because it
operates on an autoregressive language model:

\begin{center}\small
\renewcommand{\arraystretch}{1.3}
\begin{tabular}{@{}lll@{}}
\toprule
\textbf{Phase} & \textbf{Bottleneck} & \textbf{NSA's goal} \\
\midrule
Training (forward $+$ backward) & GPU FLOP throughput (compute-bound) & Reduce total FLOPs \\
Prefilling                       & GPU FLOP throughput (compute-bound) & Reduce total FLOPs \\
Autoregressive decoding          & HBM memory bandwidth (bandwidth-bound) & Reduce KV-cache reads \\
\bottomrule
\end{tabular}
\end{center}

The decode phase is bandwidth-bound for a specific reason: tokens are
generated one at a time, and each new token must load the entire KV cache
from HBM to SRAM to attend over all prior tokens.
There is no computation to amortise that load, so the bottleneck shifts
from FLOPs to memory reads.
NSA's three mechanisms (32$\times$ compression, top-$k$ block selection,
GQA-consistent loading) each reduce the volume of KV data read per step,
which is why NSA achieves 11$\times$ faster decoding.

\ours has neither of these bottlenecks.
\RF is not an autoregressive model.
It renders a full image in one forward pass: all $N$ triangle tokens are
processed simultaneously.
There is no sequential token generation, no KV cache, and no decode phase
in the language-model sense.

The bottleneck for \ours is simply $O(N^2)$ attention:

\begin{center}\small
\renewcommand{\arraystretch}{1.3}
\begin{tabular}{@{}lcc@{}}
\toprule
\textbf{$N$ (triangles)} & \textbf{Attention pairs} & \textbf{Status} \\
\midrule
4{,}096  & $\sim$16M  & Feasible \\
8{,}192  & $\sim$67M  & Memory-tight \\
32{,}768 & $\sim$1B   & OOM --- infeasible \\
\bottomrule
\end{tabular}
\end{center}

At 32k triangles the dense attention matrix does not fit in GPU memory.
This is not a latency or throughput problem; it is a \emph{feasibility}
problem.
Dense attention cannot run at that scale at all.

The optimization goal of \ours is therefore distinct from both of NSA's
goals: reduce $O(N^2) \to O(N \cdot k)$ to make large-$N$ rendering
\emph{possible}, where $k$ is the number of tokens each query actually
attends to across the three branches.
The primary constraint is peak GPU memory; compute reduction is a
byproduct, not the target.
The metric that validates success is not latency at a fixed $N$ but
whether the model can run at $N = 32{,}768$ and whether PSNR at that
scale is competitive with a dense baseline evaluated at smaller $N$.

% ============================================================
\section{GP-Sparse Addition 1: Normal-Coherent Local Masking}
\label{sec:main}

\begin{mdframed}[linecolor=black!30, backgroundcolor=gray!5]
\textbf{Status.} In \texttt{main.tex}.
File to touch: \code{sparse\_attention.py} only.
Zero new parameters.
\end{mdframed}

\subsection{Problem}

The local branch wastes attention capacity on back-facing neighbours
and degrades shadow boundary quality near silhouette edges.
The branch selects the $k_\ell = 64$ spatially nearest triangles by
Euclidean distance, but on a typical mesh 30--40\% of these neighbours
face away from the query.
Under the rendering equation, a triangle whose normal lies in the lower
hemisphere of the query contributes zero radiance because the cosine
factor $(\mathbf{n}_i \cdot \omega)^+$ clamps at zero.

\subsection{Masking Rule}

Suppress any neighbour $j$ whose normal forms more than a $120^\circ$
angle with the query normal.
The mask is defined by a hard indicator at inference and a sigmoid
soft approximation for STE gradient flow during training:
\begin{align}
  \tilde{m}_{ij} &= \sigma\!\bigl(-s\,(\mathbf{n}_i \cdot \mathbf{n}_j + \tau)\bigr),
  \label{eq:ncmask_soft_a2} \\
  m_{ij}         &= \mathbf{1}[\mathbf{n}_i \cdot \mathbf{n}_j < -\tau],
  \label{eq:ncmask_hard_a2}
\end{align}
where $\tau = 0.5$ and sharpness $s = 20$.
The forward pass uses $m_{ij}$ (hard, physically exact); the backward
pass propagates gradients through $\tilde{m}_{ij}$, so surface normals
receive gradient signal from the rendering loss (straight-through
estimator, Bengio et al.\ 2013).
Subtracting $\mathbf{m}_i^{\text{STE}} \cdot \infty$ from the attention
logits before softmax sets masked weights to zero.
If all neighbours of a query are masked (concave corner), unmask all
of them to prevent \code{softmax([-inf,\ldots,-inf])} returning NaN.

\textbf{Why $\tau$ is fixed.}
$\tau = 0.5$ is the hemisphere constraint of the rendering equation,
not a learned preference.
Allowing the optimiser to lower $\tau$ below zero would permit attention
to back-facing triangles, which violates the physics.

\begin{claimbox}
\textbf{Novelty.} This is back-face culling applied as an attention mask.
No prior sparse attention method applies a geometric visibility primitive
inside the attention pattern because none operates on triangle meshes.
The mask reduces effective neighbourhood size by 30--40\% at zero
inference overhead.
\end{claimbox}

\subsection{Inspiration}

\textbf{Back-face culling} is a standard technique in classical rasterisation
pipelines.
Any triangle whose surface normal faces away from the camera is discarded
before drawing because it contributes nothing to the visible image.
The normal-coherent mask applies the same physical reasoning inside the
attention pattern: a triangle whose normal lies in the lower hemisphere of
the query cannot exchange direct radiance with it, because the cosine term
$(\mathbf{n}_i \cdot \omega)^+$ in the rendering equation clamps to zero
for those directions.

The three-branch structure of GP-Sparse is structurally inspired by NSA
(Native Sparse Attention, Yuan et al., arXiv:2502.11089), which decomposes
attention into compressed-block, selected-block, and local-window branches.
NSA operates on one-dimensional token sequences with no geometric structure.
In GP-Sparse, block boundaries follow the scene BVH rather than sequence
order, pooling weights follow the rendering integral rather than uniform
averaging, and the local branch applies the hemisphere constraint as a hard
mask.
The normal-coherent mask has no counterpart in NSA.

\subsection{Disadvantages}

\begin{itemize}[nosep]
  \item \textbf{Translucent and subsurface-scattering materials.}
    These surfaces genuinely receive light from behind the hemisphere.
    The mask suppresses those interactions incorrectly, so the method
    degrades on wax, skin, leaves, and other translucent geometry.
  \item \textbf{Depends on correct surface normals.}
    Noisy, inconsistent, or flipped normals in raw mesh data break the
    mask silently: valid neighbours get suppressed and invalid ones are
    kept, with no error signal.
  \item \textbf{Fallback negates the benefit on concave geometry.}
    When all neighbours of a query are masked (tight corners, interior
    cavities), the guard unmasks all of them and the query reverts to
    full local attention. The most geometrically complex triangles
    receive no benefit.
  \item \textbf{Diffuse hemisphere only.}
    $\tau = 0.5$ encodes the diffuse hemisphere. For specular surfaces
    the relevant window is a narrow reflection lobe, not the full
    hemisphere, so the threshold is too coarse to mask correctly.
\end{itemize}

\subsection{Implementation (Key Steps)}

\begin{enumerate}[nosep]
  \item Locate the normal slice in the token tensor.
    Set \code{NORMAL\_SLICE = slice(0, 3)} and verify
    \code{n.norm(dim=-1).mean()}~$\approx 1.0$.
  \item Compute cosine of shape $(N, k_\ell)$:\\
    \code{cos\_nn = (n\_q[:, None] * n\_nbrs).sum(-1)}
  \item Build hard mask (forward / inference):\\
    \code{mask\_hard = cos\_nn < -tau}
  \item Build soft mask (backward only, STE):\\
    \code{mask\_soft = torch.sigmoid(-s * (cos\_nn + tau))}\\
    \code{mask\_ste = mask\_soft + (mask\_hard.float() - mask\_soft).detach()}
  \item Add all-masked guard:\\
    \code{mask\_hard = mask\_hard \& \textasciitilde mask\_hard.all(dim=-1, keepdim=True)}.
  \item Apply STE mask: \code{logits.masked\_fill(mask\_hard, float('-inf'))}.\\
    (STE: replace \code{mask\_hard} with \code{mask\_ste} when gradients are needed.)
\end{enumerate}

\begin{actionbox}
\textbf{Sanity check.}
On the Cornell box, the floor (normal $+y$) and ceiling (normal $-y$)
have dot product $-1 < -0.5$.
The ceiling should be masked when the floor is the query.
\end{actionbox}

\subsection{Ablation Plan}

Run at $N = 16{,}384$ on the Objaverse held-out split.

\begin{center}\small
\begin{tabularx}{\textwidth}{@{}lX@{}}
\toprule
Condition & Purpose \\
\midrule
\ours, no mask ($\tau = \infty$) & Baseline for this addition \\
\ours, $\tau = 0.5$ & Main result \\
$\tau \in \{0.0,\,0.5,\,0.9\}$ sweep & Inference-only; no retraining \\
Sparsity measurement (\% masked) & Report in paper; expected 25--45\% \\
\bottomrule
\end{tabularx}
\end{center}

% ============================================================
\section{GP-Sparse Addition 2: Area-Weighted Coarse Pooling}
\label{sec:supporting}

\begin{mdframed}[linecolor=black!30, backgroundcolor=gray!5]
\textbf{Status.} To be implemented.
Files to touch: \code{bvh.py} (precompute areas) and
\code{sparse\_attention.py} (replace \code{.mean(0)} with weighted sum).
Zero new parameters.
\end{mdframed}

\subsection{Problem}

Uniform average pooling dilutes bright emitters in leaves where they
sit alongside many dark triangles, misrepresenting the radiometric
content of the block.
The coarse branch compresses each BVH leaf node into a single token
by averaging all triangle tokens with equal weight.
Under the rendering equation, however, a surface patch's contribution
scales with its area: a large emitter and a small occluder are not
equal contributors, so equal weighting is physically incorrect.

\subsection{Fix: Area-Weighted Pooling}

Replace uniform averaging with an area-weighted sum:
\begin{equation}
  \hat{\mathbf{K}}_\beta
    = \frac{\sum_{j \in \mathcal{B}_\beta} A_j\,\mathbf{K}_j}
           {\sum_{j \in \mathcal{B}_\beta} A_j},
  \qquad
  \hat{\mathbf{V}}_\beta
    = \frac{\sum_{j \in \mathcal{B}_\beta} A_j\,\mathbf{V}_j}
           {\sum_{j \in \mathcal{B}_\beta} A_j},
  \label{eq:area_pool}
\end{equation}
where $A_j = \tfrac{1}{2}\|(v_1^j - v_0^j)\times(v_2^j - v_0^j)\|$
is the area of triangle $j$.

Area weighting is the leading-order approximation to solid angle:
for compact BVH leaves where triangles are spatially close and similarly
oriented, $A_j / r^2$ (the solid angle contribution) reduces to $A_j$
up to a shared scale factor.
The weights are query-independent, computed once per BVH build, and
add no parameters.

\textbf{Boundary fringe.}
Standard area-weighted pooling assigns each triangle to exactly one leaf.
A triangle at the boundary between two leaves therefore contributes nothing
to the neighbouring leaf's pooled token.
If that triangle is an emitter, the block scorer will fail to assign high
importance to the neighbouring leaf, and the selected branch will not
retrieve it.
To close this gap, the $k_{\text{fringe}} = 4$ triangles from neighbouring
leaves that are geometrically closest to each leaf's boundary triangles are
included in the area-weighted pooling with their true areas.
Fringe membership is determined at BVH build time and adds no per-forward-pass
overhead.
The fringe size of 4 is a conservative choice; ablation over
$k_{\text{fringe}} \in \{0, 2, 4, 8\}$ should be included in the
experimental section.

\begin{claimbox}
\textbf{Novelty.} With both additions in place, the coarse and local
branches are grounded in the rendering equation: the coarse branch uses
area weighting consistent with the rendering integral, and the local
branch applies the hemisphere constraint as a hard mask
(Section~\ref{sec:main}).
To our knowledge, no prior sparse attention method applies
physics-based weighting to its pooling step.
\end{claimbox}

\subsection{Implementation}

\step{1} In \code{bvh.py}, precompute areas at BVH build time:
\begin{verbatim}
e1 = verts[:, 1] - verts[:, 0]        # (N, 3)
e2 = verts[:, 2] - verts[:, 0]
areas = 0.5 * torch.cross(e1, e2, dim=-1).norm(dim=-1)  # (N,)
\end{verbatim}

\step{2} In the coarse pooling loop, replace \code{.mean(0)} with:
\begin{verbatim}
w = areas[block_idx]                   # (b,)
k_hat = (tokens[block_idx] * w[:, None]).sum(0) / w.sum()
\end{verbatim}

\step{3} Sanity check: set all non-emitter emission features to zero
in a leaf that contains one large emitter.
The pooled token's emission component should match the emitter's value
weighted by its area fraction.

\subsection{Ablation Plan}

\begin{center}\small
\begin{tabularx}{\textwidth}{@{}lX@{}}
\toprule
Condition & Purpose \\
\midrule
\ours with uniform pooling & Baseline for this addition \\
\ours with area-weighted pooling & Main result \\
PSNR on emitter-heavy scenes vs.\ general scenes & Tests the hypothesis \\
\bottomrule
\end{tabularx}
\end{center}

\subsection{Disadvantages}

\begin{itemize}[nosep]
  \item \textbf{Weights are query-independent.}
    Area weight is the same for all queries attending to a block.
    A large triangle far from the query still receives high weight even
    though it subtends a small solid angle from that distance.
    The approximation holds for compact leaves but degrades when queries
    are spread far from the block.
  \item \textbf{Minimal benefit when SAH produces homogeneous leaves.}
    Surface area heuristic splitting tends to group triangles of similar
    size into the same leaf. When leaf triangles are all similarly sized,
    area weights are nearly uniform and the weighted sum reduces to an
    ordinary average.
  \item \textbf{Does not help fragmented emitters.}
    If a large light source spans multiple BVH leaves, its signal is
    split across multiple pooled tokens regardless of weighting within
    each leaf. Area weighting within a leaf cannot recover a signal that
    is distributed across leaves.
\end{itemize}

% ============================================================
\section*{Combined Checklist}

\begin{itemize}[nosep]
  \item[$\square$] Main: implement mask and verify on Cornell box
  \item[$\square$] Main: measure sparsity on 100 held-out scenes
  \item[$\square$] Main: run $\tau$ sweep at inference time on trained checkpoint
  \item[$\square$] Supporting: precompute triangle areas in \code{bvh.py}
  \item[$\square$] Supporting: replace \code{.mean()} with area-weighted pooling
  \item[$\square$] Supporting: sanity check on emitter-heavy leaf node
  \item[$\square$] Supporting: ablation on emitter-heavy vs.\ general split
  \item[$\square$] Update contributions and coarse-branch equation in \texttt{main.tex}
\end{itemize}

\vspace{10pt}
\hrule
{\small\color{gray} Internal document. Do not distribute.}

% ============================================================
\newpage
\section*{GPU Requirements}

\subsection*{Short answer}

No, this work does not require heavy GPU resources.
A single A100-80\,GB (or A6000-48\,GB for ablations) is sufficient.

\subsection*{Why fine-tuning is enough}

The original \RF needed 8$\times$A100s for 8 days to train from scratch on
16M Objaverse renders.
This project does not do that.
Fine-tuning starts from the released checkpoint and re-initialises only
the attention projection matrices in the view-independent stage (roughly
10--15\% of the 205M parameters).
Everything else --- the Swin-L pixel decoder, the view-dependent stage,
the token embeddings --- stays frozen.

\subsection*{Why GP-Sparse makes single-GPU fine-tuning possible}

Dense attention at $N = 32{,}768$ would require approximately $64\times$
more memory than at $N = 4{,}096$, making even a single forward pass
infeasible on a multi-GPU cluster.
GP-Sparse reduces the view-independent stage from $O(N^2)$ to effectively
$O(N)$, which is what makes fine-tuning at the target scale tractable on
one GPU in the first place.

\subsection*{Estimated requirements by task}

\begin{center}\small
\renewcommand{\arraystretch}{1.3}
\begin{tabular}{@{}lll@{}}
\toprule
\textbf{Task} & \textbf{GPU} & \textbf{Approx.\ time} \\
\midrule
Stage~1 warm-up (scorer KL, frozen weights) & 1$\times$A100-80\,GB & Hours \\
Stage~2 sparse fine-tuning ($N = 16$k scenes) & 1$\times$A100-80\,GB & Days, not weeks \\
Inference / evaluation at $N = 32$k & 1$\times$A100-80\,GB & Seconds per scene \\
Ablations at $N = 16$k & 1$\times$A6000-48\,GB & Workable \\
Blender Cycles dataset renders & CPU or any GPU & Overnight batches \\
\bottomrule
\end{tabular}
\end{center}

\subsection*{The real bottleneck is data, not training}

Generating 1{,}000 scenes at $N = 16$k with Blender Cycles at 512\,spp
takes overnight GPU time.
The action steps document flags this explicitly: start rendering in Week~2,
not after implementation finishes.
Training time is not the constraint; waiting for data is.

\subsection*{How to get free compute}

Apply to all three sources immediately --- approval takes one to two weeks:

\begin{itemize}[nosep]
  \item AWS Research Credits ---
    \url{https://aws.amazon.com/research-credits}
  \item Google TPU Research Cloud ---
    \url{https://sites.research.google/trc}
  \item Microsoft Accelerate Foundation Models Research (AFMR) ---
    {\small\url{https://www.microsoft.com/en-us/research/collaboration/accelerate-foundation-models-research}}
\end{itemize}

\vspace{10pt}
\hrule
{\small\color{gray} Internal document. Do not distribute.}

\end{document}
